\unit{Squeezing, coherence and photon correlations}
\unitref{W3L2, W3L3; W6L3; Assignment 1}
\subsection{What is actually being squeezed?}
Squeezing changes \emph{field quadrature
uncertainties}. A vacuum state has equal
variance 1 in \(x,p\) with our convention.
For real \(r\), a squeezing operator can be defined as
\[
S(r)=\exp\left[\frac r2(a^2-a^{\dagger2})\right].
\]
Its Bogoliubov action is
\[
S^\dagger aS=a\cosh r-a^\dagger\sinh r.
\]
Consequently
\[
S^\dagger xS=e^{-r}x,\qquad
S^\dagger pS=e^{r}p,
\]
and the squeezed vacuum \(S(r)\ket0\) has
\[
\boxed{(\Delta x)^2=e^{-2r},\quad
(\Delta p)^2=e^{2r},\quad
\Delta x\,\Delta p=1.}
\]
The uncertainty ellipse becomes narrow in
one direction and wide in the other;
the product does \emph{not} violate Heisenberg.

\begin{center}
\begin{tikzpicture}[>=Latex,scale=.8]
\draw[->](-3,0)--(3,0)node[right]{\(x\)};
\draw[->](0,-2.5)--(0,2.5)node[above]{\(p\)};
\draw[dashed,gray,thick](0,0)circle(1.05);
\draw[accent,thick,fill=accent!10](0,0)ellipse(.55 and 1.95);
\node[below right]at(.5,-1.65){squeezed};
\node[above left,gray]at(-1,1){vacuum};
\end{tikzpicture}
\end{center}

\subsection{How can vacuum contain squeezed photons?}
The squeezed vacuum has
\(\langle a\rangle=0\), but its mean
photon number is
\[
\boxed{\langle n\rangle=\sinh^2r.}
\]
It contains even-number superposition components
produced by a parametric process.
A \emph{displaced squeezed} state has nonzero
mean field as well as anisotropic noise.
The order of \(D(\alpha)\) and \(S(r)\)
matters if a calculation asks for the
total number population, because
the two operators do not generally commute.

\subsection{Squeezing helps only the measured direction}
If a signal shifts the \(x\)-mean by
\(\delta x\), the noise-limited sensitivity
scales as \(\delta x/\Delta x\).
An \(x\)-squeezed probe reduces
\(\Delta x\) by \(e^{-r}\) and improves this
signal-to-noise ratio. But if one measures
the \(p\) quadrature instead, the noise
is \emph{larger}. The optical phase of
the local oscillator must therefore
match the squeezing axis.

\begin{worked}{How loss spoils squeezing}
Take input \(x\)-quadrature variance
\(V_{\rm in}=e^{-2r}=1/4\), and a channel
with power transmission \(\eta=0.8\).
A vacuum environment enters through the
unobserved beam-splitter port:
\[
V_{\rm out}=\eta V_{\rm in}+(1-\eta)V_{\rm vac}
=0.8(1/4)+0.2(1)=\boxed{0.40}.
\]
The output is still squeezed relative
to vacuum (\(V_{\rm out}<1\)), but
not as strongly as the input.
\end{worked}

\subsection{What does coherence measure?}
First-order coherence measures field
correlations across space or time:
\[
G^{(1)}(t,t+\tau)
=\langle a^\dagger(t)a(t+\tau)\rangle.
\]
It is related to ordinary interference
visibility and spectral linewidth.
Second-order coherence measures correlated
photon arrivals,
\[
G^{(2)}(t,t+\tau)=
\langle a^\dagger(t)a^\dagger(t+\tau)
a(t+\tau)a(t)\rangle.
\]
Normalised \(g^{(2)}(0)>1\) describes
bunching; \(g^{(2)}(0)<1\) describes
antibunching, which is nonclassical
for a stationary single-mode field
under conventional assumptions.
An experiment can reveal coherence
without being a direct measurement
of a particular optical phase.

\subsection{Pair production and heralding}
A two-mode squeezed vacuum
\(\sum_n\lambda^n\ket{n,n}\)
contains signal--idler number correlations.
If the idler is unobserved,
the signal looks thermal. If a detector
heralds exactly one idler photon,
the conditional signal is \(\ket1\).
This is possible because measurement
changes the state of the \emph{remaining}
subsystem, even if the global evolution
before detection was unitary.

\begin{checkpoint}
Squeezing concerns quadrature noise.
Bunching concerns photon arrival statistics.
First-order coherence concerns field
correlations. These are related but
not interchangeable ideas.
\end{checkpoint}
